\documentclass[12pt]{article}
\usepackage[brazil]{babel}
\addto\captionsbrazil{
  \renewcommand{\tablename}{Tabela}
  \renewcommand{\figurename}{Figura}
}
%template SBC
\usepackage{amssymb}
\usepackage{sbc-template}
\usepackage{graphicx,url}
\usepackage[utf8]{inputenc}
\usepackage{array}
\usepackage[table]{xcolor}
\usepackage{adjustbox}
\usepackage{float}
\usepackage{tikz}
\usetikzlibrary{shapes.geometric}
\usetikzlibrary{arrows.meta, positioning, shapes.geometric}
\usepackage{amsfonts}
\usepackage{multirow}
\usepackage{pgfplots}
\pgfplotsset{compat=1.18}
\usepackage{amsmath}
\usetikzlibrary{shapes,arrows}
\usepackage{booktabs} % Para formatação de tabelas profissionais
\usetikzlibrary{positioning, backgrounds, fit, matrix, shapes.geometric}
\definecolor{techblue}{RGB}{31, 78, 121}
\definecolor{techorange}{RGB}{227, 108, 10}
\definecolor{bglight}{RGB}{245, 247, 250}
\definecolor{bordergray}{RGB}{120, 120, 120}
% Pacotes para citações e links clicáveis
\usepackage{cite}               % agrupa citações consecutivas
\usepackage[hidelinks]{hyperref} % torna citações e URLs clicáveis

%tabelas
\usepackage{booktabs}
\usepackage{tabularx}
\newcolumntype{Y}{>{\raggedright\arraybackslash}X}
\usepackage{titlesec}

\sloppy
\setlength{\textfloatsep}{8pt plus 2pt minus 2pt} % reduz espaço entre figuras/tabelas e o texto
\setlength{\intextsep}{8pt plus 2pt minus 2pt}

\title{Predicting Asthma Exacerbations in Digital Health Ecosystems Using Machine Learning Frameworks}
\author{Herbert Jung Sin\inst{1}, Martony Demes da Silva\inst{1}}

\begin{document}
\address{Curso Superior de Tecnologia de Gestão de Dados  \\ Centro de Educação Aberta e a Distância (CEAD) \\ Universidade Federal do Piauí (UFPI)\\
  Teresina -- PI -- Brasil
  \email{herbert.sin@ufpi.edu.br, martony.silva@ufpi.edu.br}
}

\makeatletter
\def\inst#1{}
\def\instnum{}
\makeatletter

\maketitle

\begin{abstract}
Screening the risk of acute asthma exacerbations is critical for digital health solutions and Patient-Generated Health Data (PGHD) integration. This study evaluates an intelligent screening model designed for the \textit{Afilaxy} mobile application to classify risk profiles based on clinical criteria from the Global Initiative for Asthma (GINA). Four machine learning algorithms were tested on a dataset of asthmatic profiles: Logistic Regression, Random Forest, Gradient Boosting, and XGBoost. The XGBoost model achieved an Area Under the Curve (AUC) of 91.3\% and a Sensitivity (\textit{Recall}) of 89.2\% at the optimal decision threshold. Although these results indicate strong classification capacity on synthetic proxy data, the remaining 10.8\% false-negative rate warrants clinical caution. This work presents a proof-of-concept for intelligent risk screening, establishing a foundation for future implementation on edge devices.
\end{abstract}
     
\begin{resumo}
A triagem do risco de exacerbações asmáticas agudas é crítica para soluções de saúde digital e integração de dados gerados por pacientes (PGHD). Este trabalho avalia um modelo de triagem inteligente projetado para o aplicativo móvel \textit{Afilaxy}, visando classificar perfis de risco de crise com base nos critérios clínicos da \textit{Global Initiative for Asthma} (GINA). Quatro algoritmos de aprendizado de máquina foram testados sobre uma base de dados de perfis asmáticos: Regressão Logística, Random Forest, Gradient Boosting e XGBoost. O modelo em XGBoost obteve uma Área Sob a Curva (AUC) de 91,3\% e Sensibilidade (\textit{Recall}) de 89,2\% no limiar de decisão otimizado. Embora esses resultados indiquem forte capacidade de classificação em dados proxy sintéticos, a taxa remanescente de 10,8\% de falsos negativos exige cautela na interpretação clínica. O estudo apresenta uma prova de conceito para triagem inteligente de risco, fundamentando hipóteses para futura implementação em dispositivos de borda.
\end{resumo}

\section{Introdução} \label{sec:introducao}

A asma brônquica afeta mais de 260 milhões de pessoas no mundo, sendo responsável por 450 mil mortes anuais evitáveis \cite{ginasthma2024}. No Brasil, figura entre as principais causas de internação no SUS, gerando impacto socioeconômico por custos diretos e absenteísmo \cite{sbt2021}. Embora o manejo eficaz dependa do reconhecimento precoce de sinais inflamatórios pré-crise, o monitoramento ambulatorial exibe um \textit{vazio assistencial}: um hiato temporal entre consultas esporádicas que força decisões baseadas em relatos retrospectivos e imprecisos dos pacientes.

Como consequência desse modelo reativo, os dados coletados tardiamente podem ser afetados por vieses de memória e subestimação de sintomas leves. Preditores fundamentais, como variações acentuadas no índice de qualidade do ar (\textit{Air Quality Index} - AQI) e flutuações microclimáticas de umidade e temperatura, passam despercebidos pela triagem convencional até culminarem em crises agudas e urgências hospitalares. Para parametrizar a gravidade dessas ocorrências, as diretrizes da \textit{Global Initiative for Asthma} (GINA) categorizam o controle clínico do paciente com base nos sintomas das últimas quatro semanas (Tabela~\ref{tab:gina_criteria}).

\begin{table}[ht!]
\centering
\small
\caption{Critérios GINA para Avaliação do Controle da Asma}
\label{tab:gina_criteria}
\vskip 0.05in
\begin{tabular}{p{0.46\textwidth}p{0.48\textwidth}}
\hline
\textbf{Sintomas avaliados (últimas 4 semanas)} & \textbf{Classificação do Controle (Itens Ativos)} \\ \hline
1. Sintomas diurnos ($>2\times$/semana)          & \textbf{Controlada:} Nenhum item presente \\
2. Despertares noturnos por asma\textsuperscript{*} & \textbf{Parcialmente Controlada:} 1 a 2 itens \\
3. Uso de medicação de alívio ($>2\times$/semana) & \textbf{Não Controlada:} 3 a 4 itens \\
4. Qualquer limitação de atividades físicas      & \\ \hline
\end{tabular}
\vskip 0.05in
\begin{minipage}{\textwidth}
\footnotesize{\textit{Fonte: Adaptado de GINA (2024). *Nota: O campo de sintomas noturnos integra o cálculo do score e o treinamento, enquanto a variável de tratamento preventivo (\textit{onControllerMedication}), introduzida na v2.3.1 do aplicativo, permanece suspensa por ausência desse vetor no dataset proxy utilizado.}}
\end{minipage}
\end{table}

A expansão de dispositivos móveis viabilizou a coleta de Dados Gerados por Pacientes (\textit{Patient-Generated Health Data} - PGHD). Contudo, traduzir esse fluxo heterogêneo em alertas clínicos acionáveis é um desafio, visto que scores heurísticos falham em mapear correlações não lineares. Para mitigar isso, este trabalho introduz a modelagem preditiva no arcabouço do aplicativo \textit{Afilaxy} \cite{afilaxy2025} (ativo em Android e em beta no iOS). Visando substituir sua lógica atual de pesos fixos por modelos inteligentes locais (\textit{on-device}) sob escassez inicial de dados históricos nativos, adota-se um \textit{dataset proxy} consolidado, enriquecido algoritmicamente para emular os módulos de \textit{check-in} e sensores do aplicativo.

O objetivo geral deste artigo é avaliar a viabilidade teórica e estatística de algoritmos de aprendizado de máquina supervisionado para a classificação do risco de exacerbação asmática (definido pelo rótulo \textit{high\_risk} segundo as diretrizes GINA). Em termos pragmáticos, a "viabilidade computacional" investigada neste estudo refere-se especificamente à capacidade dos algoritmos de obter alto desempenho preditivo (priorizando a Sensibilidade) mantendo uma arquitetura de baixa complexidade paramétrica compatível com restrições de inferência móvel. A principal contribuição científica deste trabalho reside na formulação de uma abordagem de triagem integrada ao ambiente \textit{Afilaxy}, combinando a adaptação das diretrizes clínicas GINA para rotulamento automático e o estudo comparativo de modelos com foco em segurança clínica, estabelecendo os requisitos teóricos para a futura implantação em borda.

\section{Referencial Teórico} \label{sec:referencial}

O desenvolvimento de um modelo preditivo para o aplicativo \textit{Afilaxy} exige a convergência entre a caracterização clínica das exacerbações, o processamento de fluxos contínuos de dados móveis, algoritmos de classificação e métricas voltadas à segurança em saúde.

\subsection{Fisiopatologia da Asma, GINA e o Papel dos Dados de Saúde (PGHD)}
A asma é uma doença inflamatória crônica das vias aéreas inferiores, caracterizada por hiperresponsividade brônquica e limitação variável ao fluxo aéreo, cujos sintomas (sibilância, dispneia e tosse) oscilam frente a gatilhos exógenos, como infecções e variações no Índice de Qualidade do Ar (\textit{Air Quality Index} - AQI) \cite{guarnieri2014, ginasthma2024}. As exacerbações representam o agravamento desse processo. O manejo preconizado pela \textit{Global Initiative for Asthma} (GINA) enfatiza o monitoramento ativo em curtas janelas temporais para estratificar o risco de crises \cite{ginasthma2024}, justificando o uso de ferramentas digitais para automatizar essa detecção.

Nesse cenário, os Dados Gerados por Pacientes (\textit{Patient-Generated Health Data} - PGHD) coletados via saúde móvel (\textit{mHealth}) reconfiguraram o monitoramento clínico. Ao contrário dos registros transversais de Prontuários Eletrônicos (\textit{Electronic Health Records} - EHR), o PGHD capta o cotidiano do indivíduo de forma longitudinal e contínua através de sintomas autorrelatados e telemetria climática passiva (temperatura, umidade e AQI por geolocalização) \cite{demiris2019}. Essa fusão de dados provê o arcabouço informacional necessário para treinar algoritmos preditivos supervisionados \cite{tsang2022}.

\subsection{Algoritmos de Classificação Supervisionada}
Para mapear o risco de crises a partir do PGHD, quatro arquiteturas destacam-se pela robustez e eficiência computacional, cobrindo desde vieses lineares interpretáveis até abordagens de comitês (\textit{ensembles}).

\begin{description}
    \setlength{\itemsep}{1pt}
    \setlength{\parskip}{0pt}
    \setlength{\parsep}{0pt}
    \item[\textbullet~Regressão Logística:] Modelo probabilístico linear que estima a ocorrência de um evento aplicando a função sigmóide sobre uma combinação linear ponderada dos atributos de entrada \cite{hastie2009}. Embora limitada para capturar interações fortemente não lineares sem engenharia prévia de atributos, sua baixa complexidade computacional e alta interpretabilidade tornam-na um \textit{baseline} indispensável no contexto médico.
    \item[\textbullet~Random Forest:] Método baseado em aprendizado em conjunto (\textit{ensemble learning}) que constrói uma multiplicidade de árvores de decisão independentes durante o treinamento \cite{breiman2001}. O modelo mitiga o sobreajuste (\textit{overfitting}) injetando aleatoriedade na amostragem dos dados via \textit{bagging} e na seleção restrita de atributos \cite{james2021}, consolidando a predição final por votação majoritária.
    \item[\textbullet~GBM e XGBoost:] Diferente do aprendizado paralelo, as arquiteturas de \textit{boosting} operam sequencialmente. O GBM minimiza a função de perda treinando novas árvores sequenciais \cite{hastie2009}, enquanto o \textit{XGBoost} introduz penalidades de regularização estrutural ($L_1$ e $L_2$) em sua função de custo, prevenindo o sobreajuste. Sua arquitetura paralela em nível de \textit{hardware} otimiza a execução local (\textit{on-device}) no sistema \textit{Afilaxy}.
\end{description}

\subsection{Sensibilidade (Recall) como Métrica de Segurança Clínica}
A validação de modelos aplicados à medicina diagnóstica difere de problemas genéricos onde a acurácia global é priorizada. Em ambientes de saúde digital, o custo clínico de um Falso Negativo (classificar um paciente em risco iminente como estável) sobrepuja drasticamente o impacto de um Falso Positivo, visto que a omissão de um alerta impede uma intervenção preventiva oportuna \cite{altman1994}. Diante dessa assimetria, a Sensibilidade (ou \textit{Recall}) assume o papel de métrica central de desempenho, pois quantifica a proporção exata de casos de risco real que o pipeline computacional detectou com sucesso \cite{altman1994}, garantindo o patamar de segurança mandatório para ferramentas de suporte à decisão.

\section{Trabalhos Relacionados} \label{sec:relacionados}

O monitoramento contínuo da asma via soluções digitais impulsiona soluções de IA focadas no cruzamento de dados clínicos com telemetria ambiental para prever exacerbações. \citeonline{bousquet2017} demonstraram, no estudo piloto MASK, a viabilidade do autorrelato diário de sintomas via aplicativo móvel, embora limitado a heurísticas estáticas. De forma complementar, \citeonline{merchant2016} acoplaram sensores a inaladores de resgate para monitorar em tempo real o uso de medicação de alívio (SABA), observando redução desse uso em um ensaio clínico randomizado.

Quanto aos modelos classificadores, a revisão clínica de \citeonline{tsang2022} mapeou estudos que aplicaram aprendizado de máquina a dados de saúde móvel para o manejo da asma e constatou o predomínio de técnicas tradicionais de ML, alertando que a maioria dos algoritmos foi desenvolvida em bases pequenas e populações restritas. Já \citeonline{genes2018} descreveram a integração de dados autorrelatados em um aplicativo de asma ao prontuário eletrônico, evidenciando os desafios técnicos de levar PGHD à prática clínica. Neste trabalho, sustenta-se ainda que um modelo embarcado de apoio à decisão deve priorizar alta sensibilidade e manter baixa complexidade estrutural, compatível com a execução local (\textit{on-device AI}). A Tabela~\ref{tab:trabalhos_relacionados} sintetiza o panorama.

\begin{table}[ht!]
\centering
\small
\caption{Comparativo metodológico dos trabalhos relacionados}
\label{tab:trabalhos_relacionados}
\vskip 0.05in
\begin{tabularx}{\textwidth}{lXXXX}
\hline
\textbf{Autor (Ano)} & \textbf{Técnicas Base} & \textbf{Origem Dados} & \textbf{Variáveis} & \textbf{Foco Principal} \\ \hline
Bousquet (2017) & Heurísticas Clínicas & App Mobile Ativo & Sintomas diários & Viabilidade do autorrelato \\ \hline
Merchant (2016) & Ensaio Clínico Randomizado & Sensores em Inaladores & Uso de SABA & Monitoramento de medicação \\ \hline
Tsang (2022) & Revisão de ML & Estudos de mHealth & Variadas & Panorama de algoritmos \\ \hline
Genes (2018) & Relato de Caso & App Mobile + Prontuário & Sintomas autorrelatados & Integração de PGHD \\ \hline
\end{tabularx}
\vskip 0.05in
\footnotesize{\textit{Fonte: Elaborada pelos autores (2025).}}
\end{table}

Apesar dos avanços na captação de PGHD e no benchmark de classificadores, observa-se na literatura uma lacuna de engenharia: a maioria dos modelos robustos funciona em silos na nuvem, dependendo de conexão constante à internet e ignorando restrições de infraestrutura do smartphone. O diferencial deste trabalho reside no desenvolvimento de uma abordagem embarcada integrada ao sistema real do aplicativo \textit{Afilaxy}. A arquitetura proposta calibra o modelo XGBoost focado na Sensibilidade clínica ($89,2\%$ no limiar otimizado) para blindar o paciente contra falsos negativos, validando a viabilidade de conversão do artefato final para inferência leve na borda (\textit{on-device AI}), eliminando a dependência de servidores e protegendo a autonomia do dispositivo.

\section{Metodologia}

A arquitetura metodológica deste trabalho fundamenta-se na Saúde Digital Preditiva e em \textit{Machine Learning} para modelagem de risco respiratório. O desenho experimental estruturado garante o controle de sobreajuste (\textit{overfitting}) e a robustez estatística frente a dados heterogêneos (Figura~\ref{fig:fluxograma_asma}).\footnote{O pipeline e os scripts estão disponíveis publicamente em: \url{https://github.com/herb-sin/afilaxy-ml}}

\begin{figure}[ht!]
\centering
\small

\tikzset{
    block/.style={rectangle, draw=techblue, fill=white, minimum width=3cm, minimum height=1.2cm, text width=2.8cm, align=center, rounded corners=2pt, thick},
    database/.style={rectangle, draw=techblue, fill=techblue!15, minimum width=2.2cm, minimum height=1.2cm, text width=2cm, align=center, rounded corners=6pt, thick},
    layer/.style={draw=bordergray, dashed, fill=bglight, inner sep=12pt, rounded corners=4pt, thick},
    arrow/.style={thick,->,>=stealth}
}

\begin{tikzpicture}[node distance=1.1cm and 1.0cm, font=\sffamily\small]

    % CAMADA 1: INGESTÃO E ENGENHARIA (TOP)
    \node (db1) [database] {\textbf{PGHD} \\ \scriptsize{Diários / Wearables}};
    \node (p1) [block, right=of db1] {\textbf{1.1 Pré-processamento} \\ \scriptsize{Imputação \& Min-Max}};
    \node (p2) [block, right=of p1] {\textbf{1.2 Janelamento} \\ \scriptsize{Médias Móveis (7d/14d)}};
    \node (vector) [block, draw=techorange, fill=techorange!15, right=of p2, minimum width=2.2cm, text width=2cm] {\textbf{Vetor} \\ $\mathbf{X_{norm}}$};

    \draw[arrow] (db1) -- (p1);
    \draw[arrow] (p1) -- (p2);
    \draw[arrow] (p2) -- (vector);

    \begin{scope}[on background layer]
        \node (group1) [layer, fit=(db1) (p1) (p2) (vector)] {};
        \node[anchor=north west, techblue] at (group1.north west) {\textbf{1. Pipeline de Engenharia de Dados}};
    \end{scope}

    % CAMADA 2: MODELAGEM E OTIMIZAÇÃO (MIDDLE)
    \node (rf_block) [block, below=1.8cm of p1, minimum width=3.5cm, text width=3.3cm] {\textbf{2.1 Modelos de ML} \\ \scriptsize{LR, RF, GBM e XGBoost}};
    \node (grid_block) [block, right=of rf_block] {\textbf{2.2 Grid Search \& CV} \\ \scriptsize{Otimização ($K=5$)}};

    \draw[thick,->,>=stealth, techorange] (vector.south) -- ++(0,-0.6) -| (rf_block.north);
    \draw[arrow] (rf_block) -- (grid_block);

    \begin{scope}[on background layer]
        \node (group2) [layer, fit=(rf_block) (grid_block)] {};
        \node[anchor=north west, techblue] at (group2.north west) {\textbf{2. Inteligência Computacional / Classificadores}};
    \end{scope}

    % CAMADA 3: VALIDAÇÃO OPERACIONAL (BOTTOM)
    \node (val_block) [block, below=1.8cm of rf_block] {\textbf{3.1 Validação K-Fold} \\ \scriptsize{Estratificação ($K=5$)}};
    \node (metrics_block) [block, right=of val_block] {\textbf{3.2 Avaliação} \\ \scriptsize{Recall \& AUC-ROC}};
    \node (db2) [database, right=of metrics_block] {\textbf{Métricas} \\ \scriptsize{AUC / F1-Score / ROC}};

    \draw[thick,->,>=stealth, techblue] (grid_block.south) -- ++(0,-0.6) -| (val_block.north);
    \draw[arrow] (val_block) -- (metrics_block);
    \draw[arrow] (metrics_block) -- (db2);

    \begin{scope}[on background layer]
        \node (group3) [layer, fit=(val_block) (metrics_block) (db2)] {};
        \node[anchor=north west, techblue] at (group3.north west) {\textbf{3. Avaliação Estatística e Diagnóstica}};
    \end{scope}

\end{tikzpicture}
\caption{Arquitetura de dados ponta a ponta e engenharia de inteligência preditiva para monitoramento da asma crônica.}
\label{fig:fluxograma_asma}
\end{figure}

\subsection{Definição do Target e Justificativa Clínica}

A modelagem no \textit{Afilaxy} foca na predição de exacerbações agudas no intervalo entre consultas, e não no diagnóstico de asma. O rótulo original da base (\textit{Diagnosis}) foi descartado para mitigar \textit{data leakage}. Em seu lugar, estruturou-se a variável binária \textit{high\_risk} ($y \in \{0, 1\}$) via engenharia de atributos sobre o histórico de sintomas.

A rotulagem baseou-se nas diretrizes da GINA, contabilizando cinco sintomas de \textit{check-ins} diários: chiado (\textit{Wheezing}), falta de ar (\textit{ShortnessOfBreath}), aperto no peito (\textit{ChestTightness}), tosse (\textit{Coughing}) e sintomas noturnos (\textit{NighttimeSymptoms}). O escore ($\text{\textit{symptom\_score}} \in [0, 5]$) mapeia o controle da asma segundo a Equação~(\ref{eq:target}), definindo alto risco a partir de 3 indicadores:

\begin{equation}
\label{eq:target}
\textit{high\_risk} = 
\begin{cases} 
1, & \text{se } \text{\textit{symptom\_score}} \ge 3 \\
0, & \text{se } \text{\textit{symptom\_score}} < 3
\end{cases}
\end{equation}

Após a parametrização nos 2.392 registros da base \textit{proxy}, a distribuição do \textit{target} mostrou-se naturalmente balanceada, dispensando amostragem artificial: a classe \textit{high\_risk} = 1 representa \textbf{58,1\%} (1.390 amostras) e \textit{high\_risk} = 0 representa \textbf{41,9\%} (1.002 amostras). Essa proporção é consistente com a prevalência individual dos sintomas (50\%–60\%) e próxima ao modelo binomial esperado ($57,7\%$), descartando viés no desenho do rótulo.


\subsection{Descrição Detalhada das Etapas Experimentais}
A execução experimental da metodologia proposta é estruturada em quatro etapas encadeadas, que englobam desde a estruturação e pré-processamento dos dados brutos até a otimização de hiperparâmetros e validação estatística dos classificadores, conforme detalhado a seguir.

\subsubsection{Etapa 1: Ingestão e Tratamento de Dados}
A proposta inicia-se com a coleta e agregação de dados de saúde gerados por pacientes (PGHD - \textit{Patient-Generated Health Data}) em conformidade teórica com a perspectiva de Demiris et al. (2019). O referencial metodológico para o pré-processamento segue as diretrizes de integridade de dados biológicos e de telemetria digital:
\begin{itemize}
    \item \textbf{Tratamento de Omissões e Ruídos:} Registros nulos provenientes de falhas de sincronização de dispositivos vestíveis ou lapsos de preenchimento de diários digitais são tratados via imputação baseada na mediana móvel do paciente para variáveis contínuas, ou pela moda local para variáveis categóricas, prevenindo vieses de exclusão.
    \item \textbf{Normalização:} Para garantir a estabilidade computacional e a convergência dos algoritmos, aplica-se a transformação de escala linear Min-Max. Este procedimento realiza o mapeamento dos atributos numéricos originais uniformemente dentro do intervalo centrado entre zero e um, neutralizando disparidades de magnitude entre diferentes unidades de medida sem distorcer a dispersão dos dados.
\end{itemize}

\subsubsection{Etapa 2: Engenharia de Atributos (\textit{Feature Engineering})}
Para capacitar os classificadores a capturarem a dinâmica das exacerbações de asma, a base de dados bruta é transformada através da extração de três grupos de atributos preditivos:
\begin{enumerate}
    \item \textbf{Janelas Temporais de Sintomas e Função Pulmonar:} Estruturação de variáveis agregadas (médias e desvios padrões móveis de 7 e 14 dias) a partir do score de sintomas diários relatados e dos parâmetros espirométricos dinâmicos de Volume Expiratório Forçado ($FEV_1$) e Capacidade Vital Forçada ($FVC$), permitindo mapear declínios funcionais graduais à luz das recomendações globais da GINA \cite{ginasthma2024}.
    \item \textbf{Variáveis Ambientais Exógenas:} Alinhamento georreferenciado de dados macroclimáticos (temperatura, umidade relativa do ar) e níveis de poluentes particulados ($PM_{2.5}$ e $PM_{10}$)\footnote{Na versão v2.3.1 do app, apenas o índice AQI composto é armazenado no banco de dados. Embora as concentrações individuais de $PM_{2.5}$ e $PM_{10}$ sejam consultadas via API da WAQI, elas não são persistidas no histórico do paciente, exigindo um ajuste de modelo de dados antes do retreinamento.}, fundamentando-se na evidência clínica estabelecida por Guarnieri e Balmes (2014) sobre o impacto da poluição na hiperresponsividade brônquica.
    \item \textbf{Variabilidade de Variáveis Fisiológicas e do Ecossistema Mobile:} Mapeamento do desvio padrão móvel da frequência cardíaca, saturação de oxigênio ($SpO_2$) e padrões de eficiência do sono. Embora a arquitetura do aplicativo Afilaxy esteja preparada via estruturas (\textit{scaffolds}) no HealthConnect e HealthKit para integrar tais módulos, eles permanecem inativos na versão v2.3.1 (retornando nulo no modelo de dados). Por essa razão, e pela ausência desses vetores na base proxy laboratorial, essas variáveis foram omitidas do treinamento do classificador final.
\end{enumerate}

\subsubsection{Etapa 3: Modelagem Preditiva e Otimização}

O núcleo de inteligência computacional foi estruturado a partir de uma abordagem comparativa com quatro classes de algoritmos complementares: Regressão Logística (adotada como \textit{baseline}), \textit{Random Forest}, \textit{Gradient Boosting Machine} (GBM) e o \textit{XGBoost}. 

Para extrair o teto de desempenho e garantir uma comparação equitativa, as arquiteturas foram submetidas a um processo de busca exaustiva em grade (\textit{Grid Search}) integrado a um esquema de validação cruzada estratificada ($K=5$) restrito ao subconjunto de treinamento. O espaço de busca multidimensional concentrou-se nos parâmetros estruturais mais críticos de cada modelo:
\begin{itemize}
    \item \textbf{Regressão Logística:} Força da regularização inversa ($C \in \{0.01, 0.1, 1.0, 10.0\}$);
    \item \textbf{Modelos baseados em Árvores (RF, GBM e XGBoost):} Volume do comitê ($n\_estimators \in \{50, 100, 200\}$) e profundidade máxima das ramificações ($max\_depth \in \{3, 5, 7\}$);
    \item \textbf{Algoritmos de Boosting (GBM e XGBoost):} Fator de encolhimento ou taxa de aprendizado ($learning\_rate \in \{0.01, 0.05, 0.1, 0.2\}$).
\end{itemize}

Os hiperparâmetros ótimos selecionados automaticamente em cada partição foram consolidados nos artefatos finais para a etapa subsequente de validação estatística e diagnóstica.

\subsubsection{Etapa 4: Avaliação de Desempenho e Validação Estatística}
Para mitigar o risco de estimativas excessivamente otimistas e garantir a generalização do modelo para novos indivíduos, adotou-se um protocolo de validação cruzada estratificada em 5 partições ($K=5$) no conjunto de treino, seguido pela avaliação final em um conjunto de testes independente contendo 20\% dos registros totais. Devido à reestruturação da variável resposta fundamentada nas diretrizes GINA, que resultou em uma distribuição equilibrada de classes ($58,1\%$ alto risco versus $41,9\%$ baixo risco), o pipeline dispensou a aplicação de algoritmos de amostragem artificial como o SMOTE, preservando a distribuição epidemiológica real dos dados.

O desempenho preditivo da arquitetura foi quantificado analiticamente a partir da matriz de confusão, fundamentando-se nos conceitos de validação de testes diagnósticos descritos por Altman e Bland (1994). O comportamento do classificador final foi verificado por meio da extração da Sensibilidade (\textit{Recall}), Especificidade, Valor Preditivo Positivo (VPP), Valor Preditivo Negativo (VPN) e a Área Sob a Curva ROC (AUC-ROC), assegurando o equilíbrio ideal entre a segurança clínica do usuário e a especificidade operacional do aplicativo.

\subsection{Formulação da Tarefa Preditiva e Limitações Metodológicas} \label{subsec:limitacoes_metodologia}
É importante pontuar que a formulação experimental adotada neste trabalho configura uma \textit{classificação de estado/triagem de risco corrente} baseada na presença de sintomas ativos e variáveis exógenas simultâneas, e não uma \textit{predição temporal antecipada} com janela de horizonte futuro delimitada (ex.: prever uma crise com $N$ dias de antecedência). 

Adicionalmente, como o rótulo $high\_risk$ é derivado das diretrizes GINA aplicando regras sobre os sintomas autorrelatados ($Wheezing$, $ShortnessOfBreath$, etc.), esses indicadores foram \textbf{excluídos} da matriz de atributos $X$ para prevenir vazamento direto de dados (\textit{data leakage}). Em seu lugar, a variável sintética $riskScore$ — calculada a partir de exposição ambiental, função pulmonar e histórico de sintomas — introduz uma sobreposição indireta entre preditores e \textit{target}. Esse fator contribui parcialmente para os valores de AUC observados, devendo o modelo ser interpretado como um emulador da lógica de triagem GINA adaptada a aplicação \textit{Afilaxy}, e não como uma prova de aprendizado causal autônomo.

\section{Resultados e Discussão} \label{sec:resultados}

A análise descritiva da base \textit{proxy} (2.392 registros, 15 variáveis) abrange: 5 atributos clínico-demográficos (idade de 5–80 anos, gênero, etnia, tabagismo e IMC de 15,0–40,0 kg/m²); 2 espirométricos ($FEV_1 \approx 2,54$ L/s e $FVC \approx 3,00$ L/s); 4 ambientais/comportamentais (poluição, pólen, poeira e atividade física); e 4 variáveis mapeadas ao \textit{Afilaxy} (\textit{riskScore}, crise prévia, inalador de resgate e acionamentos de emergência). Para evitar \textit{data leakage}, os 5 sintomas ativos foram excluídos da matriz $X$ (Seção~\ref{subsec:limitacoes_metodologia}). Essa estrutura multifatorial simula o sistema heterogêneo do aplicativo para a modelagem da estabilidade clínica.

\subsection{Análise do Desempenho Global dos Classificadores}

Os algoritmos foram avaliados em um conjunto de testes independente com 20\% da base original. A reestruturação do \textit{target} via diretrizes GINA resultou em uma distribuição equilibrada para o aprendizado supervisionado, refletindo a realidade epidemiológica sem a necessidade de amostragem artificial.

A otimização via \textit{Grid Search} com validação cruzada no treino definiu os hiperparâmetros do modelo \textit{XGBoost}: taxa de aprendizado (\textit{learning\_rate}) de $0,20$, profundidade máxima (\textit{max\_depth}) igual a $3$ e $100$ estimadores ($n\_estimators=100$). A profundidade rasa atuou como regularizador estrutural, reduzindo o risco de sobreajuste nos dados transversais e favorecendo a capacidade de generalização do classificador.

A Tabela~\ref{tab:metricas_finais} consolida as métricas obtidas por cada arquitetura na fase de avaliação.

\begin{table}[ht!]
\centering
\small
\caption{Desempenho Comparativo dos Modelos no Conjunto de Teste}
\label{tab:metricas_finais}
\vskip 0.1in
\begin{tabularx}{\textwidth}{lXXXX}
\hline
\textbf{Modelo} & \textbf{Média CV (AUC)} & \textbf{AUC (Teste)} & \textbf{Sensibilidade (Recall)} & \textbf{F1-Score} \\ \hline
Regressão Logística & 0,942 & \textbf{0,933} & 0,867 & \textbf{0,870} \\ \hline
Random Forest       & 0,862 & 0,881 & \textbf{0,899} & 0,843 \\ \hline
Gradient Boosting   & 0,921 & 0,910 & 0,860 & 0,854 \\ \hline
\textbf{XGBoost}    & \textbf{0,922} & 0,913 & 0,867 & 0,859 \\ \hline
\end{tabularx}
\vskip 0.05in
\footnotesize{\textit{Fonte: Dados extraídos dos experimentos computacionais (2025).}}
\end{table}

Os dados demonstram que os modelos de \textit{boosting} (XGBoost e Gradient Boosting) e a Regressão Logística operaram em faixa competitiva, com AUC entre $0{,}910$ e $0{,}933$ no conjunto de teste. A Regressão Logística obteve a maior AUC individual ($0{,}933$) e o maior $F1\text{-}Score$ ($0{,}870$), enquanto o Random Forest registrou o maior Recall ($0{,}899$). O XGBoost ($\text{AUC}=0{,}913$; $F1=0{,}859$) apresentou desempenho competitivo com os demais modelos de boosting. Na ausência de testes estatísticos de hipótese formais (como Wilcoxon ou teste $t$ pareado nas dobras da validação cruzada), não se deve inferir superioridade estatística estrita de uma arquitetura sobre a outra. O XGBoost é adotado como modelo de referência para implantação (\textit{deployment}) em razão de suas propriedades de regularização estrutural ($L_1$ e $L_2$) e compatibilidade com inferência em borda (\textit{on-device AI}), aspectos críticos no sistema \textit{Afilaxy}.

\subsection{A Métrica de Sensibilidade como Critério de Segurança Clínica}
Sob a ótica de segurança em saúde digital, onde a omissão de um alerta preventivo pode levar ao agravamento do quadro respiratório, o modelo \textit{Random Forest} obteve o maior Recall no limiar padrão (\textbf{89,9\%}), enquanto XGBoost e Regressão Logística alcançaram \textbf{86,7\%}. Com ajuste do limiar de decisão (Seção~\ref{subsec:threshold}), o Recall do XGBoost pode ser elevado para \textbf{89,2\%}.

Embora essa taxa represente uma alta capacidade de detecção no conjunto testado, é indispensável reconhecer a relevância clínica da margem de erro remanescente: um \textit{Recall} de $89,2\%$ implica que \textbf{10,8\% dos indivíduos em situação de alto risco não foram identificados pelo modelo} (falsos negativos). Em uma coorte em produção, essa taxa representa pacientes vulneráveis sem o devido alerta. Portanto, o resultado deve ser interpretado como um avanço na triagem automatizada, mas que não substitui o acompanhamento médico e nem garante a eliminação absoluta de crises.

\subsection{Impacto da Otimização do Limiar de Decisão} \label{subsec:threshold}

A Figura~\ref{fig:threshold_plot} ilustra o comportamento dinâmico das métricas de validação frente à variação do limiar probabilístico do modelo XGBoost. O mapeamento visual evidencia o \textit{trade-off} clássico em saúde: a elevação do limiar aumenta a Precisão, porém provoca queda acentuada na Sensibilidade.

\begin{figure}[ht!]
\centering
\includegraphics[width=0.85\textwidth]{figura2_threshold.png}
\vskip 0.05in
\caption{Comportamento das Métricas de Validação em Função do Limiar de Decisão}
\label{fig:threshold_plot}
\vskip 0.05in
\footnotesize{\textit{Fonte: Gráfico extraído do pipeline de calibração do experimento (2025).}}
\end{figure}

Para o sistema \textit{Afilaxy}, o ponto de operação padrão em $0,50$ foi substituído pela calibração do limiar em \textbf{0,466} (sinalizado pela linha tracejada), elevando o \textit{Recall} para \textbf{89,2\%}. O limiar foi calibrado para maximizar o F1-Score, o que elevou o Recall sem degradar a precisão de forma proporcional — um compromisso alinhado ao princípio de precaução clínica, que prioriza a captura de exacerbações mesmo ao custo de um pequeno número adicional de falsos positivos.

\subsection{Análise de Relevância dos Atributos} \label{subsec:importancia}

Para conferir interpretabilidade ao modelo de referência, extraiu-se a importância de atributos (\textit{Feature Importance}) do algoritmo \textit{XGBoost}, baseada no ganho relativo médio trazido para as divisões estruturais do \textit{ensemble} (Figura~\ref{fig:feature_importance}).

\begin{figure}[ht!]
\centering
\includegraphics[width=0.85\textwidth]{figura3_feature_importance.png}
\vskip 0.05in
\caption{Importância dos Atributos no XGBoost (\textit{Dataset Proxy}).}
\label{fig:feature_importance}
\vskip 0.05in
\footnotesize{\textit{Fonte: Dados do experimento (2025).}}
\end{figure}

Essa distribuição reflete a capacidade discriminativa no dataset, \textbf{sem implicar causalidade direta ou mecanismos fisiopatológicos}.

A variável \textit{emergencyCount} (contagem de acionamentos do botão de emergência no app) emergiu como o preditor de maior peso (25,8\%), seguida por \textit{riskScore} (18,0\%) — saída do motor heurístico do \textit{Afilaxy} — e \textit{PollutionExposure} (12,3\%). Os indicadores espirométricos $FEV_1$ (7,1\%) e $PollenExposure$ (6,4\%) completam os cinco atributos mais relevantes, convergindo com o conhecimento clínico sobre gatilhos de exacerbação asmática. Notadamente, a inclusão do $riskScore$ como \textit{feature} cria dependência do modelo com a regra heurística estática, devendo ser removida em estudos de ablação futuros. Tais padrões mostram-se \textbf{compatíveis com o conhecimento clínico}, devendo ser validados prospectivamente com dados longitudinais nativos do aplicativo.

\subsection{Discussão de Viabilidade Técnica e Implantação no App Real}
A avaliação do modelo provê diretrizes para o desenvolvimento contínuo do aplicativo \textbf{Afilaxy}:

\begin{itemize}
    \item \textbf{Superação do Motor Heurístico:} A transição de regras fixas para modelos baseados em árvores permite identificar padrões não lineares entre clima e sintomas, tornando a triagem de risco mais dinâmica.
    \item \textbf{Uso de Dados Proxy e Necessidade de Validação Longitudinal:} É indispensável reforçar que os resultados numéricos foram obtidos em uma base sintética e transversal (\textit{Asthma Disease Dataset}). Em produção real, o sistema enfrentará séries temporais ruidosas, falhas de preenchimento por parte do usuário e ausência de atributos demográficos e clínicos basais (como idade, IMC e histórico de tabagismo, que atualmente não são coletados pelo perfil de usuário do app), exigindo novos testes em dados longitudinais nativos.
    \item \textbf{Hipótese de Inferência na Borda (On-Device AI):} A conversão do modelo para \textit{.tflite} ou \textit{.mlmodel} fundamenta a viabilidade de execução local, mas essa integração não está implementada na v2.3.1. Além disso, a computação das janelas de 7 e 14 dias exige um pipeline local (ou Cloud Function) para agregar o histórico do Firestore. Salienta-se que \textbf{não foram realizados testes empíricos de bateria, latência, RAM ou tamanho de arquivo}, sendo essa validação um trabalho futuro indispensável.
    \item \textbf{Dependência de Hardware para Parâmetros Espirométricos:} A utilização das variáveis de função pulmonar ($FEV_1$ e $FVC$) no modelo pressupõe a integração com espirômetros portáteis certificados via Bluetooth. Esta barreira de hardware deve ser considerada no roadmap de implantação do solução.
\end{itemize}

\section{Conclusão} \label{sec:conclusao}

Este trabalho apresentou uma prova de conceito para triagem inteligente do risco de exacerbação asmática no aplicativo \textit{Afilaxy}. O modelo \textit{XGBoost} demonstrou elevada capacidade discriminativa em dados sintéticos, alcançando AUC de 91,3\% e Sensibilidade de 89,2\% no limiar otimizado, superando abordagens puramente heurísticas no monitoramento contínuo.

Como limitações, destacam-se o viés do \textit{dataset proxy} transversal, a sobreposição funcional entre atributos e o rótulo de risco, e a ausência de testes de desempenho em hardware móvel. Destaca-se que o \textit{Afilaxy} possui arquitetura para integrar 7 variáveis adicionais (como frequência cardíaca e $SpO_2$ via HealthConnect/HealthKit), atualmente pendentes de ativação nativa. Trabalhos futuros devem focar no retreinamento dos modelos com dados longitudinais reais, testes de significância estatística, modelagem temporal preditiva e \textit{benchmarking} empírico de latência e consumo energético da inferência \textit{on-device}.

\bibliographystyle{sbc}

\bibliography{referencias}

\end{document}